\chapter{Injection or Banding of Haemorrhoids}
\index{Injection or Banding of Haemorrhoids}

\section*{Definition and Overview}
Injection and banding are outpatient treatments for internal haemorrhoids that have not settled with dietary and lifestyle measures. Rubber band ligation, the most common procedure, involves placing a small elastic band around the base of the haemorrhoid so that its blood supply is cut off and it shrinks and falls off within days. Injection sclerotherapy involves injecting a chemical into the haemorrhoid to shrink it. Both are quick, do not require general anaesthesia, and allow most people to return to normal activities immediately. They are highly effective for early and moderately advanced internal haemorrhoids, and are performed by GPs with special training, colorectal surgeons, or specialist nurses.

\section*{Epidemiology}
Haemorrhoids affect up to half of adults at some time, with prevalence peaking between ages 45 and 65. Around half of people with significant haemorrhoids require some form of treatment, and rubber band ligation is among the most commonly performed procedures, with several hundred thousand performed worldwide each year. Internal haemorrhoids are more common in people who strain at stool, have chronic constipation or diarrhoea, are overweight, are pregnant, or sit for long periods. The risk of prolapse and bleeding increases with age and with worsening of these factors.

\section*{Aetiology and Pathophysiology}
Haemorrhoids are normal cushions of blood vessels in the anal canal that help maintain continence. They become symptomatic when the supporting tissues stretch and the cushions prolapse, often because of increased pressure from straining, constipation, pregnancy, or prolonged sitting. Internal haemorrhoids are graded by degree of prolapse: first degree bleed without prolapsing; second degree prolapse on straining but reduce spontaneously; third degree require manual reduction; and fourth degree remain prolapsed. Banding and injection target first to third degree internal haemorrhoids by cutting off the blood supply (banding) or causing inflammation and scarring (injection), so that the haemorrhoid shrinks and becomes attached to the underlying tissue.

\section*{Clinical Features}
\begin{itemize}
    \item Bright red bleeding on the toilet paper or in the toilet bowl
    \item Prolapse of soft tissue from the anus, which may need to be pushed back
    \item Itching, irritation, and mucus discharge
    \item Discomfort or a feeling of incomplete emptying
    \item Pain is usually absent with internal haemorrhoids unless thrombosed or after banding
    \item After banding: a mild ache and a small amount of bleeding for a few days
    \item Bleeding after injection is usually minimal
\end{itemize}

\section*{Differential Diagnosis}
Bleeding from the back passage has other causes that must be excluded.
\begin{itemize}
    \item \textbf{Anal fissure:} a painful tear in the anal lining
    \item \textbf{Colorectal cancer:} bleeding in older people always warrants investigation
    \item \textbf{Inflammatory bowel disease:} colitis with bleeding and diarrhoea
    \item \textbf{Diverticular disease:} bleeding from the colon
    \item \textbf{Polyps:} growths that may bleed
    \item \textbf{Angiodysplasia:} abnormal blood vessels in the bowel
    \item \textbf{Thrombosed external haemorrhoid:} a painful, tender lump at the anal margin
\end{itemize}

\section*{When to Seek Medical Care}
Anyone with bleeding from the back passage should have it assessed, as it can be caused by serious conditions. People with persistent haemorrhoid symptoms that do not respond to fibre, fluids, and straining avoidance should discuss treatment options. Pain, heavy bleeding, or changes in bowel habit warrant prompt review.

\textbf{Contact your local emergency services for:}
\begin{itemize}
    \item Heavy rectal bleeding, especially with faintness, dizziness, or collapse
    \item Severe pain, fever, or a swollen, red area around the anus
    \item Inability to pass urine after treatment
    \item Sudden severe anal pain with a firm, tense lump
\end{itemize}

\section*{Diagnosis and Investigations}
\begin{itemize}
    \item \textbf{History:} bleeding pattern, bowel habit, and risk factors
    \item \textbf{Examination:} inspection and a digital rectal examination
    \item \textbf{Proctoscopy:} direct inspection of the anal canal to grade the haemorrhoids
    \item \textbf{Colonoscopy or flexible sigmoidoscopy:} recommended in older people, those with anaemia, or when cancer needs exclusion
    \item \textbf{Full blood count:} to assess anaemia from chronic bleeding
    \item \textbf{Assessment of suitability:} the grade and position of the haemorrhoids determine whether banding or injection is appropriate
\end{itemize}

\section*{Management}
\begin{itemize}
    \item \textbf{Rubber band ligation:} one to three small bands placed at the haemorrhoid base, repeated at intervals if needed
    \item \textbf{Injection sclerotherapy:} a chemical injected to shrink early haemorrhoids
    \item \textbf{Lifestyle measures:} high-fibre diet, fluids, and avoiding straining
    \item \textbf{Topical treatments:} creams and suppositories for irritation
    \item \textbf{Pain relief:} simple analgesics for the mild ache after banding
    \item \textbf{Further treatment:} repeat sessions, or surgery such as haemorrhoidectomy or stapled haemorrhoidopexy for higher-grade disease
    \item \textbf{Follow-up:} review to confirm resolution and assess any further bleeding
\end{itemize}

\section*{Complications}
\begin{itemize}
    \item Pain and bleeding after the procedure, usually mild and brief
    \item Infection, rarely leading to serious sepsis, requiring prompt treatment
    \item Difficulty passing urine, especially after banding
    \item Incomplete treatment, with recurrence and the need for further sessions
    \item Rarely, banding of a higher structure causing pelvic pain
    \item Delayed bleeding a week or two after banding
\end{itemize}

\section*{Prognosis and Outcomes}
Rubber band ligation controls bleeding and prolapse in around three-quarters of people, and many require only one session. It is more effective than injection for prolapsing haemorrhoids, while injection suits early, bleeding haemorrhoids. Most people return to work the same day. Recurrence over years is common and relates to persisting constipation and straining, so attention to bowel habit remains essential. Higher-grade or recurrent haemorrhoids may need surgery, which has higher success but a more painful recovery.

\section*{Prevention and Self-Care}
\begin{itemize}
    \item Eat a high-fibre diet with plenty of fruit, vegetables, and whole grains
    \item Drink enough fluid and stay physically active
    \item Avoid straining on the toilet and prolonged sitting
    \item Do not delay opening the bowels when the urge arises
    \item Use stool softeners if needed rather than straining
    \item Maintain a healthy weight
    \item Seek early treatment for bleeding or prolapse
    \item After banding, expect a mild ache and report heavy bleeding, fever, or inability to pass urine
\end{itemize}

\section*{Sources and Further Reading}
The following reputable sources provide further information for healthcare students and patients seeking reliable, current advice about injection and banding treatment of haemorrhoids.
\begin{itemize}
    \item Healthdirect Australia. \textit{Haemorrhoids and treatment.} https://www.healthdirect.gov.au/
    \item Gastroenterological Society of Australia. \textit{Haemorrhoid treatment information.} https://www.gesa.org.au/
    \item Royal Australasian College of Surgeons. \textit{Haemorrhoid surgery information.} https://www.surgeons.org/
    \item NHS. \textit{Haemorrhoids treatment.} https://www.nhs.uk/conditions/haemorrhoids/
    \item Mayo Clinic. \textit{Haemorrhoid treatments.} https://www.mayoclinic.org/
    \item MSD Manual. \textit{Haemorrhoids.} https://www.msdmanuals.com/
\end{itemize}
